\chapter{Suites numériques}
\thispagestyle{fancy}

\begin{bhistoire}
	Dès l'Antiquité, Archimède de Syracuse (-287 ; -212), met en oeuvre une procédure itérative pour trouver une approximation du nombre $\pi$. Il encadre le cercle par des polygones inscrits et circonscrits possédant un nombre de côtés de plus en plus grand. Par ce procédé,
	Archimède donne naissance, sans le savoir, à la notion de suite numérique.\\
	Vers la fin du XVII$^\text{e}$ siècle, des méthodes semblables sont utilisées pour résoudre des équations de façon approchée pour des problèmes de longueurs, d'aires, ...\\
	Un formalisme plus rigoureux de la notion de suite n'apparaitra qu'au début du XIX$^\text{e}$ siècle avec le mathématicien français Augustin Louis Cauchy (1789 ;1857).
\end{bhistoire}


\section{généralités sur les suites numériques}

\subsection{Définition d'une suite numérique}

\begin{bexemplet}{Exemple d'introduction}
	On considère une liste de nombres formée par tous les nombres impairs rangés dans l'ordre croissant : $1,3,5,7, \ldots$\\
	On note $\left(u_{n}\right)$ l'ensemble des éléments de cette suite de nombres tel que :\\
	$u_{0}=1, u_{1}=3, u_{2}=5, u_{3}=7, \ldots$\\
	On a ainsi défini une suite numérique.\\
	On peut lui associer une fonction définie sur $\mathbb{N}$ par 
	$$u: \mathbb{N} \rightarrow \mathbb{R}$$
	$$\phantom{......................}n \mapsto u(n)=u_{n}$$
	
\end{bexemplet}

\begin{bdefi}
	\begin{itemDef}
		\item Une \grasInDef{suite numérique} $\left(u_{n}\right)$ est une liste ordonnée de nombres réels telle qu'à tout entier $n$ on associe un nombre réel noté $u_{n}$.
		\item $u_{n}$ est appelé \grasInDef{le terme de rang $\boldsymbol{n}$} de cette suite (ou d'indice $\left.n\right)$.
	\end{itemDef}
\end{bdefi}

\subsection{Suite définie par une formule explicite}

\begin{bexemples}
	\begin{itemEx}
		\item Pour tout $n$ de $\mathbb{N}$, on donne : $u_{n}=2 n$ qui définit la suite des nombres pairs. Les premiers termes de cette suite sont donc :\\
		$u_{0}=2 \times 0=0$\\
		$u_{1}=2 \times 1=2$\\
		$u_{2}=2 \times 2=4$\\
		$u_{3}=2 \times 3=6$
		\item Pour tout $n$ de $\mathbb{N}$, on donne $: v_{n}=3 n^{2}-1$ Les premiers termes de cette suite sont donc :\\
		$v_{0}=3 \times 0^{2}-1=-1$\\
		$v_{l}=3 \times 1^{2}-1=2$\\
		$v_{2}=3 \times 2^{2}-1=11$\\
		$v_{3}=3 \times 3^{2}-1=26$
	\end{itemEx}
\end{bexemples}
\begin{bdefi}
	Lorsqu'on génère une suite par une formule explicite, chaque terme de la suite est exprimé en fonction de $n$ et indépendamment des termes précédents.
\end{bdefi}

\subsection{Suite définie par une relation de récurrence}
\begin{bexemples}
	\begin{itemEx}
		\item On définit la suite $\left(u_{n}\right)$ par: $u_{0}=5$ et chaque terme de la suite est le triple de son précédent. Les premiers termes de cette suite sont donc :\\
		$u_{0}=5$\\
		$u_{1}=3 \times u_{0}=3 \times 5=15$\\
		$u_{2}=3 \times u_{1}=3 \times 15=45$\\
		De façon générale, on peut noter $: u_{n+1}=3 u_{n}$
		\item On définit la suite $\left(v_{n}\right)$ par : $v_{0}=3$ et pour tout $n$ de $\mathbb{N}, v_{n+1}=4 v_{n}-6$ Les premiers termes de cette suite sont donc :\\
		$v_{0}=3$\\
		$v_{1}=4 v_{0}-6=4 \times 3-6=6$\\
		$v_{2}=4 v_{1}-6=4 \times 6-6=18$\\
		$v_{3}=4 v_{2}-6=4 \times 18-6=66$\\
		Contrairement à une suite définie par une formule explicite, il n'est pas possible, dans l'état, de calculer par exemple $v_{13}$ sans connaître $v_{12}$.\\
		\ \\
		Cependant il est possible d'écrire un algorithme avec Python :
		\begin{center}
			\includegraphics[scale=0.4]{1.08.00-Img01}
		\end{center}
		Ou sur une calculatrice : 
		\begin{center}
			\includegraphics[scale=0.4]{1.08.00-Img02}
		\end{center}
	\end{itemEx}
\end{bexemples}
\begin{bdefi}
	Lorsqu'on génère une suite par une relation de récurrence, chaque terme de la suite s'obtient à partir de son terme précédent.Il faut connaître au moins un terme de la suite pour calculer les autres termes.
\end{bdefi}

\begin{brmq}
	Le mot \emph{récurrence} vient du latin \emph{recurrere} qui signifie "revenir en arrière".
\end{brmq}

\subsection{Représentation graphique d'une suite}

Dans un repère du plan, on représente une suite par un nuage de points de coordonnées $\left(n ; u_{n}\right) .$
\begin{bexemple}
	Pour tout $n$ de $\mathbb{N}$, on donne $: u_{n}=\dfrac{n^{2}}{2}-3$.\\
	On construit le tableau de valeurs avec les premiers termes de la suite :
	\begin{tabular}{|Xc|Xc|Xc|Xc|Xc|Xc|Xc|Xc|Xc|Xc|}
		\hline$n$ & 0 & 1 & 2 & 3 & 4 & 5 & 6 & 7 & 8 \\
		\hline$u_{n}$ & $-3$ & $-2,5$ & $-1$ & 1,5 & 5 & 9,5 & 15 & 21,5 & 29 \\
		\hline
	\end{tabular}\\
	Il est possible d'obtenir un nuage de points à l'aide d'un logiciel.
	\begin{center}
		\includegraphics[scale=0.4]{1.08.00-Img03}
	\end{center}
	
\end{bexemple}


\section{Suites arithmétiques}

\subsection{Définition}

\begin{bexemple}
	Considérons une suite numérique $U_n$ où la différence entre un terme et son précédent reste constante et égale à 5.\\
	Si le premier terme est égal à 3, les premiers termes successifs sont :\\
	$U_0 = 3$,\\
	$U_1 = 8$,\\
	$U_2 = 13$,\\
	$U_3 = 18$.\\
	Une telle suite est appelée une suite arithmétique de raison 5 et de premier terme 3.
	La suite est donc définie par : $U_{n+1}=U_n+5$ et $U_0=3$.
\end{bexemple}

\begin{bdefi}
	Une suite $U_n$ est une \grasInDef{suite arithmétique} s'il existe un nombre r tel que pour tout entier $n$, on a : $U_{n+1}=U_n+r$.
	Le nombre $r$ est appelé \grasInDef{raison} de la suite.
\end{bdefi}
\ \\
\begin{bmethodet}{Démontrer si une suite est arithmétique}
	1) La suite $U_n$ définie par : $U_n=7-9n$ est-elle arithmétique ?\\
	2) La suite $V_n$ définie par : $V_n=n^2+3$ est-elle arithmétique ?\\
	\\
	1) $U_{n+1}-U_n=7-9(n+1)-7+9n=7-9n-9-7+9n=-9$.\\
	La différence entre un terme et son précédent reste constante et égale à –9.\\
	$U_n$ est une suite arithmétique de raison $–9$.\\
	
	2) $V_{n+1}-V_n=(n+1)^2+3-n^2-3=n^2+2n+1+3-n^2-3=2n+1$.\\
	La différence entre un terme et son précédent ne reste pas constante.\\
	$V_n$ n'est pas une suite arithmétique.
\end{bmethodet}

\begin{bprop}
	$U_n$ est une suite arithmétique de raison $r$ et de premier terme $U_0$.\\
	Pour tout entier naturel $n$, on a : $U_n=U_0+nr$.
\end{bprop}

\begin{bdemo}
	La suite arithmétique $U_n$ de raison $r$ et de premier terme $U_0$ vérifie la relation $U_{n+1}=U_n+r$.\\
	En calculant les premiers termes :\\
	$U_1=u_0+r$\\ 
	$U_2=u_1+r$\\
	$U_3=u_2+r$\\
	…\\
	$U_n=U_{n-1}+r$\\
	
	En additionnant membre à membre ces n égalités, on obtient :\\
	$U_1+U_2+U_3+...+U_n=U_0+U_1+U_2+...+U_{n-1}+n\times r$\\
	Soit, en retranchant aux deux membres les termes identiques :
	$U_n=U_0+nr$
\end{bdemo}

\begin{bmethodet}{Déterminer la raison et le premier terme d'une suite arithmétique}
	Considérons la suite arithmétique $U_n$ tel que $U_5=7$ et $U_9=19$.\\
	1) Déterminer la raison et le premier terme de la suite $U_n$.\\
	2) Exprimer $U_n$ en fonction de $n$.\\ 
	\\
	1) Les termes de la suite sont de la forme $U_n=U_0+nr$\\
	Ainsi 	$U_5=U_0+5r=7$  et $U_9=U_0+9r=19$.\\
	En soustrayant membre à membre, on obtient :\\ $U_0+5r-U_0-9r=7-19$\\
	Soit : $5r-9r=7-19$ donc $r=3$.
	
	Comme $U_0+5r=7$, on a : $U_0+5\times 3=7$ et donc : $U_0=-8$.\\
	\\
	2) $U_n=U_0+nr$  soit  $U_n=-8+n\times 3$ ou encore $U_n=3n-8$
\end{bmethodet}

\subsection{Représentation graphique}

Les points de la représentation graphique d'une suite arithmétique sont alignés.

\begin{bexemple}
	On a représenté ci-dessous la suite de raison $–0,5$ et de premier terme $4$.
	\begin{center}
		\includegraphics[scale=1]{01-07-01}
	\end{center}
\end{bexemple}

\section{Somme des termes consécutifs: suite arithmétique}

\begin{bprop}
	$n$ est un entier naturel non nul alors on a : $1+2+3+...+n= \dfrac{n(n+1)}{2}$
\end{bprop}

\begin{brmq}
	Il s'agit de la somme des $n$ premiers termes d'une suite arithmétique de raison $1$ et de premier terme $1$
\end{brmq}

\begin{bdemo}
	\begin{tabular}{ c c c c c c c c c c c c }			
		& 1 & + & 2 & + & 3 & + &... & + &n-1 & + & n \\
		+ & n & + & n-1 & + & n-2 & + & ...& + & 2 & + & 1 \\
		\hline
		& (n+1) & + & (n+1) & + & (n+1) & + & ... & + & (n+1) & + & (n+1)\\
	\end{tabular} \\
	\\
	Donc : $2\times (1+2+3+...+(n-1)+n)=n(n+1)$\\
	Et donc : $1+2+3+...+(n-1)+n=\dfrac{n(n+1)}{2}$
\end{bdemo}

\begin{bmethodet}{Calculer la somme des termes d'une suite arithmétique}
	Calculer les sommes $S_1$ et $S_2$ suivantes :\\ 
	$S_1=1+2+3+...+348$\ \ \ \ \ \ \ \ \ \ \ \ \ \ \ \ \ \ \ \ \ \ \ $S_2=33+36+39+...+267$
	\\
	\\
	$S_1=1+2+3+...+348=\dfrac{348\times 349}{2}=60726$\\
	\\
	\\
	$S_2=33+36+39+...+267=3\times(11+12+13+...+89)$\\
	$=3\times((1+2+3+...+89)-(1+2+3+...+10))$\\
	$=3\times \left(\dfrac{89\times 90}{2}-\dfrac{10\times11}{2}\right)=11\ 850$
\end{bmethodet}

\section{Suites géométriques}

\subsection{Définition}

\begin{bexemple}
	Considérons une suite numérique $U_n$ où le rapport entre un terme et son précédent reste constant et égale à 2.\\
	Si le premier terme est égal à 5, les premiers termes successifs sont :\\
	$U_0 = 5$,\ \ \ 
	$U_1 = 10$,\ \ \ 
	$U_2 = 20$,\ \ \ 
	$U_3 = 40$.\\
	Une telle suite est appelée une suite géométrique de raison $2$ et de premier terme $5$.\\
	La suite est donc définie par : $U_{n+1}=2\times U_n$ et $U_0=5$. 
\end{bexemple}

\begin{bdefi}
	Une suite $U_n$ est une \grasInDef{suite géométrique} s'il existe un nombre q tel que pour tout entier $n$, on a : $U_{n+1}=q\times U_n$.
	Le nombre $q$ est appelé \grasInDef{raison} de la suite.
\end{bdefi}

\begin{bmethodet}{Démontrer si une suite est géométrique}
	La suite $U_n$ définie par : $U_n=3\times 5^n$  est-elle géométrique ?
	\ \\
	$\dfrac{U_{n+1}}{U_n} =\dfrac{3\times 5^{n+1}}{3\times 5^n}=\dfrac{5^{n+1}}{5^n} =5^{n+1-n}=5$\\
	Le rapport entre un terme et son précédent reste constant et égale à $5$.\\
	$U_n$ est une suite géométrique de raison $5$ et de premier terme $U_0=3\times 5^0=3$.
\end{bmethodet}

\begin{bexemplet}{exemple concret}
	On place un capital de \EUR{500} sur un compte dont les intérêts annuels s'élèvent à $4\%$.
	Chaque année, le capital est multiplié par $1,04$.
	Ce capital suit une progression géométrique de raison $1,04$.\\
	
	On a ainsi :\\ 
	$U_1=1,04\times 500=520$\\		
	$U_2=1,04\times 520=540,80$\\
	$U_3=1,04\times 540,80=562,432$\\	
	De manière générale : $U_{n+1}=1,04\times U_n$ avec $U_0=500$\\ 
	On peut également exprimer $U_n$ en fonction de $n$ : $U_n= 500\times1,04^n$. 
\end{bexemplet}

\begin{propriete}
	$U_n$ est une suite géométrique de raison $q$ et de premier terme $U_0$.\\
	Pour tout entier naturel $n$, on a : $U_n= U_0\times q^n$.
\end{propriete}

\begin{bdemo}
	La suite géométrique $U_n$ de raison $q$ et de premier terme $U_0$ vérifie la relation
	$U_{n+1}=q\times U_n$.
	\begin{itemDem}
		\item Si $q$ ou $U_0$ est nul, alors tous les termes de la suite sont nuls. La démonstration est évidente dans ce cas.
		\item Dans la suite, on suppose donc que $q$ et $U_0$ sont non nuls. Dans ce cas, tous les termes de la suite sont non nuls.
	\end{itemDem}
	En calculant les premiers termes :\\ 
	$U_1=q\times U_0$\\
	$U_2=q\times U_1$\\
	$U_3=q\times U_2$\\
	…\\
	$U_n=q\times U_{n-1}$\\
	\\
	En multipliant membre à membre ces $n$ égalités, on obtient :
	$u_1\times U_2\times U_3\times …\times U_n=U_0\times U_1\times U_2\times …\times u_{n-1}\times q^n$.\\
	Comme les termes de la suite sont non nuls, on peut diviser aux deux membres les facteurs identiques : $U_n=U_0\times q^n$
\end{bdemo}

\begin{bmethodet}{Déterminer la raison et le premier terme d'une suite géométrique}
	Considérons la suite géométrique $U_n$ tel que $U_4=8$ et $U_7=512$.\\
	Déterminer la raison et le premier terme de la suite $U_n$.\\
	\\
	\ \\
	Les termes de la suite sont de la forme $U_n=U_0\times q^n$.\\
	Donc : $U_4=U_0\times q^4=8$ et $U_7=U_0\times q^7=512$.\\
	Ainsi : $\dfrac{U_7}{U_4}  = \dfrac{U_0\times q^7}{U_0\times q^4}  =q^3$ et $\dfrac{U_7}{U_4}=\dfrac{512}{8}=64$ donc $q^3=64$.\\
	On utilise la fonction racine troisième de la calculatrice pour trouver le nombre qui élevé au cube donne 64.\\
	Ainsi $q=^3\sqrt{64}=4$\\
	Comme $U_0 \times q^4=8$, on a : $U_0\times 4^4=8$ et donc : $U_0= \dfrac{1}{32}$.
\end{bmethodet}

\subsection{Représentation graphique}

\begin{bexemple}
	La suite géométrique $U_n$ définie par $U_n=-4\times2^n$ est décroissante car le premier terme est négatif et la raison est supérieure à 1
	\begin{center}
		\includegraphics[scale=1]{01-07-03.png}
	\end{center}
	Remarque : Si la raison q est négative alors la suite géométrique n'est pas monotone.
\end{bexemple}
% \begin{center}
	%     \includegraphics[scale=0.8]{01-07-04.png}
	% \end{center}

\section{Somme des termes consécutifs d'une suite géométrique}

\begin{bprop}
	$n$ est un entier naturel non nul et $q$ un réel différent de 1 alors on a :\\ 
	$1+q+q^2+q^3+...+q^n= \dfrac{1-q^{n+1}}{1-q}$
\end{bprop}

\begin{brmq}
	Il s'agit de la somme des $n+1$ premiers termes d'une suite géométrique de raison $q$ et de premier terme $1$.
\end{brmq}

\begin{bdemo}
	$S=1+q+q^2+q^3+...+q^n$\\
	$q\times S=q+q^2+q^3+q^4+...+q^(n+1)$\\
	\\
	Ainsi :\\
	$S-q\times S=(1+q+q^2+q^3+...+q^n )-(q+q^2+q^3+q^4+...+q^{n+1})$\\
	$S-q\times S=1-q^(n+1)$\\
	$S\times (1-q)=1-q^{n+1}$\\
	$S=\dfrac{1-q^(n+1)}{1-q}$
\end{bdemo}

\begin{bmethodet}{Calculer la somme des termes d'une suite arithmétique}
	Calculer la somme $S$ suivante : 
	$S=1+3+3^2+3^3+...+3^{13}$\\
	\\
	$S=1+3+3^2+3^3+...+3^{13}=\dfrac{1-3^{14}}{1-3}=2\  391\  484$
	
\end{bmethodet}




